\section{Governance, Licensing, and Standards Bodies}
\label{sec:governance}

Governance is the third sub-axis of the rubric's composite
Axis 10 (\S\ref{sec:rubric:axis10}) and the dimension on which
the adoption pattern of \S\ref{sec:adoption:synthesis}
diverges most clearly from the technical rigour of the
rubric. The adoption density of A2A and MCP, despite their
modest scores on the security and authorisation axes, is
substantially explained by their institutional positioning
under the Linux Foundation Agentic AI Foundation~\cite{lf_aaif};
the limited adoption of ACNBP and LOKA, despite their
technical strength, is substantially explained by the
absence of a formal SDO standing behind either
specification. Section~\ref{sec:governance} unpacks this
institutional layer: who controls each protocol's
specification, under what license its reference
implementations ship, and what observability primitives
the specifications themselves require.

Table~\ref{tab:governance-license} consolidates the
composite-Axis-10 readings from \S\ref{sec:rubric:axis10}
into a single comparison. Section~\ref{sec:governance:foundations}
treats the open-source-foundation cluster.
Section~\ref{sec:governance:sdos} treats the activity at
formal standards-development organisations.
Section~\ref{sec:governance:tradeoff} considers the
consortium-versus-open-governance trade-off as a question
that the field has not yet decided.
Section~\ref{sec:governance:observability} treats the
observability sub-field of Axis 10 separately, because the
specification of audit-and-observability primitives is the
operational dimension on which the rubric matrix renders the
sharpest split.
Section~\ref{sec:governance:survival} closes with what
governance predicts about which protocols will still be
relevant by the end of the next standardisation cycle.

% Generated by scripts/gen_table5.py: do NOT hand-edit.
% Source: data/protocol-scores.yaml (axis_10 block).

\begin{table}[t]
  \centering
  \caption{Governance, licensing, and observability: the composite Axis 10 cell of Table~\ref{tab:rubric-matrix} expanded into its three sub-fields. Steward and license are transcribed, not scored; only observability is ordinal, on the 0 to 3 anchors of \S\ref{sec:rubric:axis10}. The rows are the eight rubric-scored protocols; the five commerce-settlement protocols are not rubric scored (\S\ref{sec:rubric:limits}) and appear in Table~\ref{tab:commerce-stack}.}
  \label{tab:governance-license}
  \footnotesize
  \begin{tabular}{@{}p{0.11\textwidth}p{0.42\textwidth}p{0.24\textwidth}c@{}}
    \toprule
    Protocol & Steward & License & Obs.\ score \\
    \midrule
    A2A & Linux Foundation & Apache 2.0 & 2 \\
    MCP & Linux Foundation AAIF & MIT-equivalent & 1 \\
    ANP & community / W3C-aligned & Apache 2.0 & 1 \\
    ACP & LF AI \& Data (archived) & Apache 2.0 & 2 \\
    AGNTCY & Linux Foundation & Apache 2.0 & 3 \\
    Coral & Coral Protocol team & CC-BY-4.0 (paper) & 2 \\
    LOKA & academic authors (CMU Tepper / BIT Sindri); no formal SDO & CC-BY-NC-SA 4.0 & 3 \\
    ACNBP & individual authors (DistributedApps.ai, Cisco, OWASP/AWS, Intuit); no formal SDO & CC-BY 4.0 & 3 \\
    \bottomrule
  \end{tabular}
\end{table}


\subsection{Open-source foundations and licenses}
\label{sec:governance:foundations}

The Linux Foundation Agentic AI Foundation (AAIF) is the
single most consequential governance institution in the
inter-agent protocol landscape at the cutoff. Founded in
December 2025 with founding members AWS, Anthropic, Block,
Bloomberg, Cloudflare, Google, Microsoft, and
OpenAI~\cite{lf_aaif}, AAIF holds stewardship of MCP, A2A,
and (via the AGNTCY collective) AGNTCY/SLIM. The
foundation's structural choice is to host protocols under
permissive open-source licenses (Apache 2.0 for A2A and
AGNTCY; an MIT-equivalent license for MCP) and to govern
their evolution through working groups that any AAIF
member can participate in.

\paragraph{The license argument.}
Apache 2.0 and MIT-equivalent licenses have the same
operational effect for the purposes of the protocols
the survey covers: anyone can implement an A2A or MCP
client or server without licensing fees or copyleft
obligations. The argument for these licenses, well-rehearsed
in the broader open-source community, is that
network-effect protocols benefit from maximum permissiveness
because any restriction shrinks the addressable
implementer population. The argument against (a strong
copyleft license, in the GPL sense, that would require
derived works to remain open) is that the contemporary
protocols already have proprietary implementations that
nobody wants to break; a copyleft license at this stage
would fragment the ecosystem rather than open it.

\paragraph{The AAIF working-group model.}
AAIF organises its substantive work in protocol-specific
and cross-cutting working groups (covering security,
identity, interoperability, and commerce) per its
public charter~\cite{lf_aaif}. The cadence at which the
working groups meet and publish revisions is faster than
W3C and IETF conventions: revision drafts have shipped on
sub-annual schedules since founding. The model is broadly
similar to the W3C and IETF working-group conventions
otherwise (public minutes, member participation, draft
specifications that mature toward stable status). It is
modulated by AAIF's smaller membership scale and faster
iteration cadence. The advantage is that specifications
can move quickly to address the empirical security gaps
the literature catalogues (\S\ref{sec:security:open}); the
disadvantage is that specifications can also break
compatibility quickly, which has begun to affect
implementers whose deployments straddle revisions.

\paragraph{Non-AAIF open-source projects.}
Coral, LOKA, ACNBP, and ANP each occupy distinct
positions outside the AAIF umbrella.
Coral Protocol Labs holds direct stewardship of the
Coral specification under an MIT license; the
governance model is a single-vendor open-source project
with a community-input mechanism but no formal working
group. LOKA's specification is authored by academic
researchers (CMU Tepper and BIT Sindri) and licensed
under CC-BY-NC-SA 4.0. That license explicitly
forbids commercial use of the specification text without
a separate agreement. This is unusual for protocol
specifications. It has been a recurring point of
discussion in the protocol-implementation community.
ACNBP's authorship is similarly informal (individual
contributors from DistributedApps.ai, Cisco, OWASP/AWS,
and Intuit) under a more permissive CC-BY 4.0. ANP is
the closest of this cluster to formal governance,
trending toward W3C through the AI Agent Protocol
Community Group~\cite{w3c_agent_protocol_cg}.

\subsection{Standards bodies}
\label{sec:governance:sdos}

Formal standards-development organisations (SDOs)
operate on slower timescales than open-source
foundations but produce outputs with stronger institutional
weight. Three SDO efforts are visible at the cutoff.

\paragraph{W3C AI Agent Protocol Community Group.}
The W3C AI Agent Protocol Community Group, established in
May 2025~\cite{w3c_agent_protocol_cg}, is the most
advanced SDO activity. The Community Group's charter
includes identity, capability descriptors, and
interoperability across heterogeneous agent
frameworks. ANP's DID-and-VC mechanism is closely
aligned with the W3C's existing DID and Verifiable
Credentials Recommendations, and the Community Group is
the venue most likely to elevate ANP-style DID-anchored
identity to W3C Recommendation status within the next
two revision cycles. The structural cost is W3C's
multi-year recommendation timeline, which is slower than
the protocols' own quarterly revision cadence.

\paragraph{IETF activity.}
The IETF has not chartered a working group specifically
for inter-agent protocols at the cutoff, but several
adjacent working groups address dimensions the protocols
need. OAuth Working Group activity on the
Agentic JWT pattern~\cite{agentic_jwt} is the most
directly relevant; secdispatch discussion of
authorisation-propagation primitives~\cite{authz_propagation}
is in the early stage that typically precedes a formal
chartering. The IETF is the institutional venue most
likely to address the authorisation-propagation gap
(O1 of \S\ref{sec:security:open}) on a multi-year
timeline.

\paragraph{ISO/IEC activity.}
ISO/IEC JTC 1/SC 42 (Artificial Intelligence) has
working items on AI agent interoperability that
overlap the inter-agent protocol space, primarily on
the governance and risk-management dimensions
(TRiSM-aligned~\cite{trism}). The ISO timeline is the
slowest of the three SDOs treated here; the visible
output at the cutoff is study reports rather than
draft standards, and the relevance to practitioners is
indirect. Where ISO/IEC outputs matter for the
inter-agent protocol landscape is in the regulatory
direction: government procurement and
sector-specific regulators increasingly defer to
ISO/IEC standards for compliance frameworks, which
gives ISO/IEC outputs a downstream effect on adoption
that exceeds their direct technical
contribution.

\subsection{Consortium versus open-governance trade-offs}
\label{sec:governance:tradeoff}

The governance question that the field has not yet
decided is the trade-off between consortium-led
governance (AAIF's model) and open standards-body
governance (W3C, IETF, ISO). The two models produce
different trajectories that the survey can characterise
without resolving.

\paragraph{The consortium argument.}
AAIF's structural advantages are speed and member
focus. Sub-annual revision cycles let the foundation
respond to empirical security findings on a timeline
matched to the threat-modelling literature's own
cadence. The founding-member set is bounded: AWS,
Anthropic, Block, Bloomberg, Cloudflare, Google,
Microsoft, and OpenAI~\cite{lf_aaif}. That bound makes
substantive disagreement tractable. A working-group meeting can reach
decisions in a small number of sessions because the
participants share enough context to argue
productively. The argument is that the contemporary
protocols are too immature for the multi-year SDO
timelines to make sense; consortium governance is the
right institutional match for the protocols' actual
maturity.

\paragraph{The open-governance argument.}
The W3C and IETF model produces specifications with
broader institutional buy-in and longer durability. A
W3C Recommendation carries weight that an AAIF working
group's output does not yet command. This holds both in
regulatory compliance contexts (where W3C outputs are
recognised in procurement specifications) and in
academic publication contexts (where W3C-derived
specifications attract citation in ways that
foundation outputs do not). The argument is that
inter-agent protocols are a generational infrastructure
layer comparable to HTTP, TLS, or DNS, and that
generational infrastructure deserves the durability
that only the slow SDO process can provide.

\paragraph{The hybrid pattern.}
The empirical observation at the cutoff is that the
two models are coexisting rather than competing. AAIF
ships the protocols at their working maturity;
W3C/IETF processes select primitives from those
working specifications and elevate them to formal
standards on a longer timeline. The W3C Community
Group's relationship to ANP is the clearest example:
ANP ships as a working specification under AAIF-style
governance; the W3C Community Group is the venue where
its DID and VC primitives are being prepared for
elevation to W3C Recommendation status. The pattern
is the inverse of the historical SDO-first model
(write the standard, then implement) and matches the
contemporary internet-infrastructure pattern
(implement, ship, then standardise).

\paragraph{The license-fragmentation risk.}
A residual risk that the consortium-versus-SDO
analysis does not capture is license fragmentation.
LOKA's CC-BY-NC-SA 4.0 license is the most-visible
example: the specification is publicly readable but
not freely re-implementable in a commercial context
without separate licensing. ACNBP's CC-BY 4.0 license
is permissive but inappropriate for protocol
specifications. CC-BY was designed for prose works,
not normative technical text. Its lack of patent
provisions has been cited in the
protocol-implementation community as a concern.
A2A, MCP, and ANP's Apache 2.0 license (with explicit
patent grants) is the cleanest license for
inter-agent protocol specifications at the cutoff;
the field would benefit from a convention that all
publicly-published specifications adopt this licensing
choice or its IETF Trust equivalent.

\subsection{Observability as the third sub-axis}
\label{sec:governance:observability}

The observability sub-field of Axis 10
(\S\ref{sec:rubric:axis10}) measures the protocol's
specification of audit, tracing, and monitoring
primitives. Table~\ref{tab:governance-license} records
the per-protocol scores; this subsection unpacks what
the scores mean operationally and why the field
distributes the way it does.

\paragraph{Score 0 baseline (unspecified).}
None of the surveyed protocols scores 0 on
observability. Every specification mentions logging
or tracing in at least a non-normative recommendation.
The baseline matters as a reference point: a protocol
that did not address observability at all would fail
production-readiness review at most enterprise
deployments.

\paragraph{Score 1 (emerging hooks).}
MCP and ANP sit at level 1. Both specifications
acknowledge that implementations should log and trace
events. Neither fixes the trace format, the
audit-record schema, or the tamper-evidence properties
of the log. ANP is the closer of the two: it requires
non-repudiation logs but standardizes no format for
reading them. The operational consequence is that
observability is a per-deployment concern rather than a
portable protocol property. Two MCP servers exposing
the same tool can log the invocation differently, and
nothing in the specification makes the two records
comparable.

\paragraph{Score 2 (trace fields specified).}
A2A, ACP, and Coral sit at level 2. The cluster is
reached by two different routes. A2A and ACP specify how
trace data is carried but not that it be trustworthy:
A2A through its SSE event stream, ACP through
OpenTelemetry integration with default exporters plus
\texttt{TrajectoryMetadata} and
\texttt{CitationMetadata} carried in the wire format
itself. ACP is the strongest cell in the group for that
reason, because putting trace metadata in the wire
format makes it portable by construction rather than by
convention. Coral arrives from the other side: \S 8
mandates a blockchain audit trail, but only for payments
and team contracts, and specifies no trace or metric
format around it. What unites the three is the missing
half. None of them specifies audit-log integrity or
non-repudiation over the whole interaction, so the audit
chain remains alterable at the operator's discretion.
The portability gain is also narrower than the level
suggests. A2A fixes an event stream but not an
audit-record schema, so an A2A monitoring system
designed for Salesforce Agentforce will not work
unmodified against ServiceNow's A2A implementation.

\paragraph{Score 3 (immutable audit + observability).}
ACNBP, LOKA, and AGNTCY sit at level 3, and they are the
three weakest-evidenced cells on this sub-field for three
different reasons. ACNBP is the cleanest: \S IV.H of the
spec~\cite{spec_acnbp} normatively requires
cryptographically protected audit logs, comprehensive
event logging, privacy-preserving logging, and
integration with external audit and compliance systems.
LOKA's level 3 is aspirational. Section 5.1 specifies
immutable audit trails, real-time ethical auditing, and
a blockchain ledger recording every action and its
reasoning, and none of it is implemented at the
cutoff~\cite{spec_loka}. AGNTCY's level 3 rests on
shipped tooling rather than on a specified property: a
dedicated observability sub-project with an
OpenTelemetry SDK, a telemetry hub, and an
observe-and-eval working group~\cite{spec_agntcy}. That
is the most usable observability stack in the surveyed
set. It is also the one cell where the anchor's first
clause is not satisfied, because OpenTelemetry
instrumentation carries trace and metric data without
making the resulting log tamper-evident, and no AGNTCY
specification text requires an audit-integrity property.
We score 3 on the strength of the second clause and flag
the cell as contestable
(\S\ref{sec:comparison:derivation}); a strict conjunctive
reading of the anchor puts AGNTCY at 2.

\paragraph{The implication for cross-tenant deployments.}
The observability sub-field is the most operationally
consequential of the three Axis 10 sub-fields for
production deployments. Multi-tenant deployments
(O4 of \S\ref{sec:security:open}) require audit
chains that survive tenant-policy disputes, and
audit chains at level 1 do not satisfy this
requirement. The SAGA TEE-anchored gateway proposal
(\S\ref{sec:security:tee}) is the institutional
response to this gap; the protocol-level response is
to elevate observability to level 3 with normative
audit-integrity properties, which only ACNBP and
AGNTCY currently do.

\subsection{What governance predicts about protocol survival}
\label{sec:governance:survival}

Three predictions about protocol survival follow from the
governance analysis. This survey is not meant to make
forecasts. However, the institutional patterns are stable
enough to support predictions on the survey's
six-to-nine-month review horizon.

\paragraph{AAIF protocols will dominate enterprise
adoption.}
A2A, MCP, and (probably) AGNTCY will continue to
dominate enterprise SaaS adoption because the AAIF
governance model matches enterprise procurement
preferences: named foundation steward, named founding
members from the enterprise vendor universe, permissive
license, monthly working-group cadence. The technical
limitations the rubric records will be addressed within
AAIF working groups on a timeline that does not require
adopters to wait for SDO outputs. The protocols may
remain weaker than ACNBP on threat coverage and weaker
than ANP on identity, but the gap will narrow as
working groups address specific gaps.

\paragraph{W3C-elevated ANP primitives will become
the identity baseline.}
The W3C Community Group's elevation of ANP-style DID
and VC primitives to W3C Recommendation status is the
single most consequential institutional event the
survey can foresee for the field. Once elevated, the
identity-layer baseline shifts from OAuth-bearer to
DID-and-VC, and A2A, MCP, and AGNTCY will adopt the
elevated primitives through their own working groups.
The timeline is uncertain. The W3C Process
Document~\cite{w3c_process_doc} specifies a multi-year
Working-Draft / Candidate-Recommendation /
Proposed-Recommendation / Recommendation pipeline. That
pipeline historically takes between eighteen months and
three years to traverse. But the direction is clear.

\paragraph{LOKA and ACNBP will remain reference
designs rather than deployed protocols.}
Neither protocol has an institutional home. LOKA rests on
academic authorship, ACNBP on individual authorship, and
neither sits under a foundation or SDO. The two protocols
will continue to be cited as the technically-rigorous
reference designs against which the deployed protocols
are measured. But they will not themselves accumulate
production deployment. The pattern is the
inter-agent analogue of CSP, Pi-calculus, and the
process-algebra family in the
concurrent-programming-language space: technically
foundational, persistently cited, never widely
deployed.

Scope compounds the absent-steward problem, and the two
should not be conflated. The argument of this paragraph is
Cisco's, made when we contacted it about this survey, and
it changed a prediction rather than a description. ACNBP
specifies discovery, identity
management, communication, and trust in one document. That
breadth is what earns it the highest threat-coverage count
in the surveyed set, and it is also what makes it hard to
adopt: an implementer cannot take the capability-binding
ceremony without also taking the naming service and the
trust model. Neither the Web nor the Internet arrived as a
single architecture. Both are collections of narrow
standards that compose, and each was adopted one piece at a
time. The prediction above therefore has a corollary that
is actionable rather than merely pessimistic. ACNBP's
mechanisms are likelier to reach production factored out
than whole, carried into A2A or MCP as separate extensions
rather than adopted as an architecture. The delegation
chain of \S\ref{sec:security:authz} is the natural first
candidate, because it is the piece with the clearest
standalone value and the fewest dependencies on the rest of
the specification.

\paragraph{The composite Axis 10 lesson.}
The composite Axis 10 reading
(\S\ref{sec:rubric:axis10}) deliberately resists
aggregation into a single number, and the governance
analysis confirms why: the steward, license, and
observability sub-fields each predict different
deployment outcomes, and combining them into a
single score would obscure the predictions that the
three sub-fields make jointly. Institutional governance,
permissive licensing, and high observability together
are necessary for production adoption at enterprise
scale; any one of them missing is a material gap that
the others cannot fully compensate for. The
field's hybrid AAIF-plus-SDO pattern
(\S\ref{sec:governance:tradeoff}) is the
institutional architecture best matched to producing
all three.
